% Data og metode

\begin{frame}{Data: A-ordningen}
\begin{itemize}\setlength{\itemsep}{4pt}
  \item Norges obligatoriske rapporteringssystem for arbeidsgivere; hele den formelle sysselsettingen.
  \item Dekning: om lag 3,1 millioner arbeidsforhold per måned, januar 2021 til april 2026.
  \item Utfall: \textbf{lønnstakerforhold med utbetalt lønn}, aktive arbeidsforhold med positiv kontantlønn i måneden.
  \item Variabler: firesifret STYRK-08-yrkeskode, alder (månedspresisjon), sektor, kontantlønn, stillingsprosent, nyansettelsesindikator.
  \item Tiårs aldersgrupper: 21--30 (tidlig karriere), 31--40, 41--50, 51--60 (senior).
  \item \textbf{Kun privat sektor}: samme kvintil rommer annet arbeid i offentlig sektor.
  \item Sesongjustert (faste månedsfaktorer, 2021--2024); indeksert til oktober 2022 = 1, måneden før ChatGPT ble lansert.
\end{itemize}
\end{frame}

% ---------------------------------------------------------------
\begin{frame}{Eksponeringsmål og kvintiler}
\begin{itemize}\setlength{\itemsep}{6pt}
  \item Hovedmål: GPT-4-eksponeringsskåren $\beta$ fra Eloundou mfl.\ (2024); dekker 397 av 407 STYRK-08-koder (97,5\,\% av sysselsettingen).
  \item Sekundært: Handa mfl.\ (2025), observert Claude-bruk delt i \textbf{automatisering} og \textbf{augmentering}; dekker 352 koder (86,5\,\%).
  \item Kodene rangeres etter eksponering og deles i fem kvintiler på \emph{yrkesfordelingen}.
  \item Hver kode teller én gang, uavhengig av sysselsettingsandel. Inndelingen ligger fast gjennom hele perioden.
\end{itemize}

\vspace{0.4em}
\footnotesize
Kontekst: Norge er nummer tre i verden i KI-bruk i befolkningen (Microsoft 2026); andelen foretak som bruker KI økte fra 11\,\% i 2021 til 30\,\% i 2025 (SSB).
\end{frame}
